\title{Large Language Models Can Automatically Engineer Features for Few-Shot Tabular Learning}

\begin{abstract}
Large Language Models (LLMs) have demonstrated substantial potential in solving complex and novel reasoning tasks, suggesting strong applicability to tabular learning—a crucial component for numerous practical domains. This study introduces a new in-context learning approach called \textsf{FeatLLM}, which harnesses LLMs specifically as feature constructors, transforming the input dataset into a form that is highly effective for making tabular predictions. The newly constructed features then inform a simple downstream prediction model (for example, linear regression), facilitating robust few-shot learning performance. Importantly, \textsf{FeatLLM} utilizes only this lightweight predictive model based on the created features during inference. Distinct from existing LLM-based methodologies, \textsf{FeatLLM} avoids the overhead of querying the LLM for every inference sample. It requires only API access to LLMs, sidestepping limitations related to prompt lengths. Empirical results over diverse tabular benchmarks across different fields indicate that \textsf{FeatLLM} formulates superior decision rules and delivers a consistent performance gain (averaging 10\%) over competing frameworks such as TabLLM and STUNT.
\end{abstract}

\section{Introduction}
Large Language Models (LLMs) have rapidly advanced, exhibiting exceptional capability in addressing previously unencountered tasks by generating appropriate responses, independent of further fine-tuning for specific tasks~\cite{creswell2022selection,imani2023mathprompter,taylor2022galactica}. Their pre-training on massive textual datasets equips LLMs with broad prior knowledge, obviating the need for human domain expertise across numerous settings~\cite{Genomes2021,singhal2023large,wilcox2003role}. This has enabled transformative impacts on task automation, notably in dialogue evaluation~\cite{finch2023leveraging} and automated program repair~\cite{fan2023automated}. By embedding instructions and limited examples within their prompts, LLMs discern the underlying context and focus on salient features or cases~\cite{brown2020language,zhao2023survey}, showcasing an aptitude for data interpretation that positions them as capable feature engineers.

Recent explorations have extended LLMs’ generalization and reasoning powers to the problem of learning from tabular data. For instance, domain-specific insights—such as correlations between age and disease susceptibility—can be integrated into prediction systems through LLM-based pipelines. Prevailing methods present both tasks and attribute details in natural language, convert the tabular data to a suitable serialized form, and subsequently pass this structured input to the LLM~\cite{dinh2022lift,wang2023anypredict}. This process may involve in-context learning~\cite{nam2023semi} or efficient parameter adaptation~\cite{dinh2022lift,hegselmann2023tabllm}. Empirically, LLM-provided priors enable enhanced accuracy, surpassing classical tabular learners especially in data-scarce situations~\cite{hegselmann2023tabllm}.

Current approaches to tabular learning with large language models (LLMs) present several significant drawbacks. For tasks requiring end-to-end prediction, the computational burden is substantial, as each data instance necessitates at least one LLM inference. The majority of these strategies depend on fine-tuning the LLM to attain satisfactory performance, limiting their utility to models whose training infrastructure or tuning APIs are accessible. This poses a considerable issue, especially since many state-of-the-art LLMs have restricted access, permitting only inference through APIs~\cite{anil2023palm,openai2023gpt4}. Moreover, existing methods are often ill-suited for handling extensive prompts~\cite{liu2023lost}; as the dimensionality of tabular data grows and the serialized prompt lengthens, both computational feasibility and model performance deteriorate. These factors collectively hinder the deployment of LLM-based solutions for tabular data in practical scenarios.

In contrast, our method reconceptualizes the role of LLMs, employing them as instruments for feature engineering rather than direct prediction. This allows for more effective exploitation of their prior knowledge. We put forth a new in-context learning paradigm, \textsf{FeatLLM}, intended to elucidate the "criteria" informing LLM decisions. For example, in a medical prediction task, the LLM can derive and articulate explicit rules associating feature values with disease classification. These identified rules serve to construct novel features that may supplant original variables in downstream analysis. This repositioning delegates feature creation to the LLM, facilitating the use of simpler models subsequently—once salient features are generated, intricate learning architectures are no longer necessary, thus reducing inference delays. The approach harnesses the LLM’s ability to learn from context, enabling feature extraction purely through prompt engineering with exemplar cases, sidestepping any requirement for model fine-tuning. Additionally, techniques inspired by ensemble learning methods~\cite{breiman2001random}, such as feature bagging, are utilized to address issues related to overly lengthy prompts.

\textsf{FeatLLM} decomposes the process of generating novel features from input prompts into two distinct sub-tasks: (1) comprehending the underlying problem and deducing the association between candidate features and targets; and (2) leveraging insights gained as well as a handful of illustrative examples to formulate discriminative rules capable of separating the classes. Through this sequenced reasoning methodology, the LLMs are guided to disregard irrelevant features during rule creation. The derived rules are subsequently implemented on the data (interpreted as code snippets), resulting in the conversion of samples into binary features that signal conformity with each rule. To assess class probabilities based on these synthesized features, a simple model is employed. Furthermore, robustness is bolstered through the use of ensemble techniques that iteratively produce features via feature bagging. By tuning the feature count and sample selection during bagging, the issue of overly lengthy prompts is addressed. \textsf{FeatLLM} operates without a need for model training and incurs minimal inference overhead, thereby making LLMs practical for a broad range of application domains.

We conduct experiments on 13 diverse tabular datasets within settings characterized by limited samples, demonstrating \textsf{FeatLLM}'s effectiveness and reliability. Our findings indicate that LLMs are adept at integrating both external prior information and dataset-derived knowledge, enabling the extraction of salient features. The proposed framework consistently surpasses state-of-the-art few-shot learning competitors across multiple configurations. The implementation is available through an anonymized GitHub repository at~\texttt{https://github.com/Sungwon-Han/FeatLLM}.

\section{Related Work}

\subsection{Few-Shot Learning with Tabular Data} The rise of few-shot learning methodologies has enabled the development of robust, transferable representations from datasets with limited labeling requirements~\cite{chen2018closer}. Originally concentrated on visual data, few-shot techniques have proven adaptable to a spectrum of modalities—including text, audio, and, most recently, tabular information~\cite{majumder2022few,nam2023stunt,schick2021exploiting}. Learning effectively from a scant number of labeled instances often leads models to pick up false or coincidental patterns~\cite{bartlett2021deep}. To counteract this, contemporary research has adopted supplementary data sources or infused domain knowledge to guide model biases. Strategies include leveraging extensive unlabeled data collections with careful data augmentation to facilitate semi-supervised learning~\cite{ucar2021subtab,yoon2020vime}, or utilizing unsupervised pre-training that is agnostic to task for optimal initialization~\cite{somepalli2021saint}. These pre-training objectives frequently involve data corruption or masking followed by reconstruction tasks~\cite{arik2021tabnet,majmundar2022met}, or deploying contrastive learning through data augmentation~\cite{bahri2022scarf,chen2020simple}. Despite these advancements, the diversity and complexity inherent in tabular data hinder the application of uniform augmentation strategies, limiting their success, especially in few-shot scenarios~\cite{nam2023stunt}. 

More recent advancements explore training models across an assortment of tabular datasets beyond those directly associated with the target task, employing transfer learning approaches~\cite{zhang2023generative,levin2022transfer,zhu2023xtab}. For example, TransTab~\cite{wang2022transtab} utilizes transformer architectures to capture semantic relationships among columns via column names rather than table entries alone. This comprehensive representation learned from disparate tables and column semantics enables models to perform few-shot or even zero-shot prediction on new tasks. Similarly, TabPFN~\cite{hollmann2022tabpfn} synthesizes data from a range of distributions for pre-training Bayesian neural networks, facilitating inference with both training and testing samples from the target task without further model adaptation. Such prior knowledge acquired from varied datasets has elevated performance beyond conventional tree-based methods, particularly in few-shot learning contexts.

\subsection{Language-Interfaced Tabular Learning} Integrating large language models (LLMs)~\cite{anil2023palm,openai2023gpt4} represents another effective method for harnessing prior knowledge in tabular learning. Due to their training on massive and heterogeneous text corpora, LLMs exhibit considerable generalization abilities across diverse tasks and disciplines~\cite{brown2020language}. Recent research efforts have focused on utilizing LLMs' rich accumulated knowledge by converting tabular information into textual formats and incorporating these serialized representations within LLM prompts~\cite{dinh2022lift,hegselmann2023tabllm,wang2023anypredict}. Input prompts containing table data, feature explanations, and task details enable the models to deduce connections between tasks and features for inference purposes. The development in this area has seen the use of parameter-efficient training mechanisms such as LoRA~\cite{hu2021lora} or IA3~\cite{liu2022few}, as well as the application of in-context learning by including several annotated examples directly in prompts, without retraining~\cite{dinh2022lift,hegselmann2023tabllm,nam2023semi,slack2023tablet}.

While the strategies discussed above demonstrate substantial progress in enhancing few-shot learning on tabular datasets, the necessity of routing every inference through an LLM presents obstacles for practical adoption in industrial environments. This research seeks to break away from the traditional paradigm where every sample’s prediction is treated as a black-box LLM task. The focus here is to engage the LLM in explicitly deducing class-specific conditions or rules. \textsf{FeatLLM} converts these rules into executable code, which subsequently generates additional features. This allows future predictions to occur independently of the LLM, and exploits in-context learning to extract rules using both prior domain knowledge and a handful of observed samples.

\section{Methods}
The goal of \textsf{FeatLLM} is to obtain ``rules"—criteria distinguishing each answer class—from a limited set of inputs, as illustrated in Figure~\ref{fig:main_model}. An LLM processes the sample inputs to generate these rules, which are then transformed into binary features that indicate, for each rule, whether a specific sample fulfills the rule or not. These binary features are assembled and utilized in a dot product with trainable non-negative weights, yielding an estimate for the probability that the sample belongs to each class. The training phase tunes the significance of each feature and reveals its contribution to class prediction (refer to Section~\ref{Sec:training}). This procedure is iterated through bagging and results are merged with an ensemble approach. The following sections elaborate on the problem formulation and provide comprehensive descriptions for each step in the methodology.

\vspace*{-0.12in}
\paragraph{Problem formulation.} Consider a tabular dataset comprised of $N$ labeled instances, denoted as $\mathcal{D} = \{(\mathbf{x}^i, \mathbf{y}^i)\}_{i=1}^{N}$. Each entry in $\mathcal{D}$ is represented by $d$ input features—where $\mathbf{x}^i$ is a vector of dimension $d$—and is accompanied by descriptive feature names in natural language, $F = \{f_j\}_{j=1}^{d}$, with corresponding definitions provided separately. Within a classification context, each label $\mathbf{y}^i$ belongs to a set of predefined categories. For $k$-shot learning scenarios, a subset of $k$ samples ($k < N$) are randomly chosen to train the model.

\subsection{Prompt Design for Extracting Rules}
\label{Sec:prompt_design}
In order to facilitate rule extraction in a manner resembling expert human reasoning, the prompt structure encourages a step-by-step approach typical of tabular learning by humans. The prompt is organized into three key sections (illustrated in Fig.~\ref{fig:prompt_for_rule} and detailed in Appendix~\ref{sec:example_rule}):

\vspace*{-0.12in}
\paragraph{Basic information description.} The opening section delivers the fundamental details required for addressing the task (corresponds to the orange segment in Fig.~\ref{fig:prompt_for_rule}). It encompasses a succinct task statement including label data, feature explanations, and illustrative example cases. The task statement is framed as a direct query (for instance, ``Does this patient's myocardial infarction complications data indicate chronic heart failure? Yes or no?"). Additional task examples appear in Appendix~\ref{sec:task_desc}. Feature descriptions specify data types and, when appropriate, definitions. Categorical features may have their possible categories enumerated. For example demonstrations, a small set of training instances are expressed in textual form, accompanied by their correct labels, using a serialization approach as outlined in previous works~\cite{dinh2022lift,hegselmann2023tabllm}:

\begin{align}
&\text{Serialize}(\mathbf{x}^i,\mathbf{y}^i, F) = \nonumber \\ 
&``\ f_1\ \text{is}\ \mathbf{x}^i_1.\ ... \ f_d\  \text{is}\  \mathbf{x}^i_d.\ \ \text{Answer:}\  \mathbf{y}^i\ ”
\end{align}

\vspace*{-0.12in}
\paragraph{Reasoning instruction.} To improve the reasoning abilities of the LLM, we introduce reasoning instructions that serve as explicit guidance (illustrated as blue text in Fig.~\ref{fig:prompt_for_rule}). These instructions start with an opening sentence modeled after the chain-of-thought methodology~\cite{wei2022chain}. The inference process is segmented into two distinct phases. Initially, the LLM is prompted to reason about the causal associations or inclinations between task-specific features and the overall task description, doing so independently from any example demonstrations and relying solely on its general or intuitive understanding. This enables the LLM to maximize its use of prior knowledge regarding the task at hand. In the subsequent phase, the LLM synthesizes information from step one along with example demonstrations to derive rule sets for each class. This bifurcated process is designed to help the model ignore extraneous correlations from non-essential data columns and better prioritize key features. To strike a balance between not extracting an insufficient number of rules (which could lead to underfitting) and avoiding an overly large prompt, we constrain the extraction to a fixed quantity of rules per class, specifically 10\footnote{Further discussion regarding rule quantity is available in Section~\ref{sec:analysis}.}. Moreover, during rule creation, the LLM refers to the feature type specified within the foundational information description, ensuring proper rule-feature compatibility.

\vspace*{-0.12in}
\paragraph{Response instruction.} For efficient parsing and subsequent usage of generated rules, we provide structured instructions that direct the LLM on how to format its outputs (see yellow text in Fig.~\ref{fig:prompt_for_rule}). These prompts detail both the prescribed response format for each class and the appropriate rule construction format (i.e., format for the rule), empowering the LLM to pick a format contingent on feature types. In the absence of further specification, the LLM has flexibility to connect multiple rules through logical connectives like AND and OR. Illustrative examples of the formatted results are included in Appendix~\ref{sec:outcome-rule}.

\subsection{Inferring Class Likelihood via Rules}
\label{Sec:training}

\paragraph{Rule parsing for feature extraction.} The rules established in the prior step are employed to construct new binary features. For each class, these features capture whether the corresponding rules are satisfied by a given sample, resulting in values of ‘0’ or ‘1’. These features serve as indicators for estimating a sample’s membership likelihood for each class. Nonetheless, because the rules are formulated in natural language by the LLM, it becomes necessary to develop a parsing protocol that allows these features to be automatically generated. Despite imposing format constraints on both the responses and rules during their creation to facilitate easier parsing, the LLM’s outputs may still display unpredictable syntactic variations~\cite{kovavc2023large}. For example, expressions such as ``$>$” or ``$<$” might be replaced by verbal forms like ``greater than” or ``less than”, complicating the parsing process.

To overcome the difficulties posed by inconsistent text formats, rather than devising intricate parsing algorithms, we employ the LLM itself as a parser. The LLM demonstrates proficiency in interpreting semantic intent even with syntax variability, and is capable of constructing concise code from instructions—leveraging its exposure to code data in pre-training. To take advantage of these capabilities, prompts containing the relevant function names, descriptions of input and output, as well as the derived rules, are provided to the LLM. Figures~\ref{fig:prompt_for_parsing} and~\ref{sec:example_parsing}-\ref{sec:example_parse_outcome} in Appendix illustrate prompt examples and resultant outputs. The code produced by the LLM is then executed via Python’s \texttt{exec()} function, enabling transformation of the dataset.

\vspace*{-0.12in}
\paragraph{Class likelihood estimation.} With binary features stemming from rule satisfaction per class, one intuitive way to assess how likely a sample belongs to a particular class is to tally the number of satisfied rules for each class—that is, sum up the binary features corresponding to each class. However, because not all rules and features contribute equally, it is important to learn their relative weight from training data. We approach this weighting using standard machine learning (ML) algorithms, applying them to the binary feature vectors of each class. For instance, a bias-free linear model could be adopted. Let $\mathbf{z}_k^i \in \{0, 1\}^R$ represent the binary features derived from sample $\mathbf{x}^i$ for class $k$, where $R$ denotes the count of rules per class and $c$ indicates the number of classes. The probability vector $\mathbf{p}^i$ for each class is then evaluated as:

\begin{align}
\text{logit}_k^i &= \max(\mathbf{w}_k, 0) \cdot \mathbf{z}_k^i, \nonumber \\
\mathbf{p}^i &= \text{Softmax}({[\text{logit}_{1}^i, ..., \text{logit}_{c}^i]}). \label{eq:infer_prob}
\end{align}

In this context, $\mathbf{w}_k$ denotes the learnable parameters in the linear model, encoding the significance of each feature. By mapping these weights into a positive domain, we guarantee that the features produced by the LLM do not adversely affect class probabilities during the training phase. The resulting logit vector undergoes transformation into class probabilities through the application of the softmax function. Training is accomplished by minimizing the cross-entropy loss, with $\mathbf{w}_k$ updated using a few-shot sample set.

\vspace*{-0.12in}
\paragraph{Ensembling with bagging.} The procedure is iteratively applied to construct several inference models, enabling final predictions to be formulated through an ensemble approach. With $T$ sets of learned weights $\{\mathbf{w}[t]\}_{t=1}^T$, the ensemble predicts by averaging class probabilities $\mathbf{p}^i$ across weights $\mathbf{w}[t]$, as detailed in Eq.~\ref{eq:infer_prob}.

Randomness is injected into rule generation for the ensemble by setting a suitably elevated temperature during LLM inference. The ordering of few-shot examples is also randomized, acknowledging prior findings that the demonstration order influences LLM responses~\cite{lu2022fantastically}\footnote{Appendix~\ref{sec:diversity} provides an analysis of rule diversity}. Bagging is used to further enhance ensemble diversity and prediction reliability by sampling a subset of features or instances for each trial, which becomes particularly beneficial as the number of features or training samples grows. The ensemble approach confers two key benefits. First, inaccuracies produced by the LLM in some trials can be offset by correctly inferred rules in others, bolstering overall system robustness. This is related to the self-consistency phenomenon in LLMs~\cite{wang2022self}, whereby recurrent rules across trials are more likely to be reliable. Second, the ensemble mechanism helps mitigate constraints imposed by the LLM's prompt length. As tabular datasets expand beyond permissible prompt size, bagging offers a way to trim input size while maintaining strong model performance. While a solitary rule set may fail to capture all feature relationships, ensembles comprised of various weak learners from multiple trials effectively reduce the impact of missing feature or instance coverage in individual models.

\section{Experiments}
We assess \textsf{FeatLLM} across a variety of tabular datasets and undertake several analytic procedures to explore the operational characteristics of our proposed framework. Supplementary experiments—including alternative LLM selection, qualitative studies, and further investigative results—are detailed in the Appendix due to limitations of main text space. 

\subsection{Performance Evaluation}
\paragraph{Datasets.}
Our study leverages 13 distinct datasets for binary or multi-class classification, including: (1) Adult~\cite{asuncion2007uci}, which focuses on predicting whether an individual's income surpasses \$50,000 per year; (2) Bank~\cite{moro2014data}, predicting the likelihood that a client subscribes to a term deposit; (3) Blood~\cite{yeh2009knowledge}, forecasting the probability of donor return; (4) Car~\cite{kadra2021well}, evaluating the quality classification of vehicles; (5) Communities~\cite{misc_communities_and_crime_183}, estimating the crime rate within a geographic area; (6) Credit-g~\cite{kadra2021well}, judging credit risk as either good or bad; (7) Diabetes\footnote{\texttt{kaggle.com/datasets/uciml/pima-indians-diabetes-database}}, for diagnosing diabetes; (8) Heart\footnote{\texttt{kaggle.com/datasets/fedesoriano/heart-failure-prediction}}, targeting detection of coronary artery disease; and (9) Myocardial~\cite{misc_myocardial_infarction_complications_579}, identifying patients with chronic heart failure.

The evaluation extends to datasets either recently released or artificially generated, which are underrepresented in existing literature and excluded from LLM pre-training corpora: (10) Cultivars, predicting the grain yield category for soybean cultivars\footnote{\texttt{archive.ics.uci.edu/dataset/913}}; (11) NHANES, classifying the age group of individuals based on record data\footnote{\texttt{archive.ics.uci.edu/dataset/887}}; (12) Sequence-type, determining the classification of a sequence among arithmetic, geometric, Fibonacci, or Collatz categories; (13) Solution-mix, deciding if a mixture's concentration from four constituent solutions exceeds 0.5. The first two were recently incorporated into the UCI Machine Learning Repository post-September 2023, guaranteeing their absence from the gpt-3.5-turbo-0613 LLM API pre-training dataset. The last two are synthetic and thus inaccessible to LLM memory, supporting unbiased evaluation.

Dataset characteristics—including size and complexity—are specified in Table~\ref{Tab:basic-desc}.
All datasets include explicit attribute names and descriptions. Notably, datasets such as `Communities' and `Myocardial' comprise more than 100 features, increasing difficulty for LLM processing due to its restricted context window.

\vspace*{-0.12in}
\paragraph{Baselines.}
Nine baseline models are incorporated for comparative study. The initial trio consists of standard tabular machine learning algorithms: (1) Logistic regression (LogReg), (2) XGBoost~\cite{chen2016xgboost}, and (3) RandomForest~\cite{ho1995random}. The subsequent pair operate with unlabeled data: (4) SCARF~\cite{bahri2022scarf}, and (5) STUNT~\cite{nam2023stunt}. It is important to highlight the practical limitations in acquiring substantial unlabeled datasets for real-world implementations. Another contemporary method is (6) TabPFN~\cite{hollmann2022tabpfn}, which leverages a model pretrained on synthetically generated data with various distributions. The remaining baselines are LLM-based approaches: (7) In-context learning (In-context)~\cite{wei2022emergent}, (8) TABLET~\cite{slack2023tablet}, and (9) TabLLM~\cite{hegselmann2023tabllm}. In-context learning relies on embedding a few annotated samples directly in the prompt, omitting parameter adjustment. TABLET improves inference by appending rule-based and prototype hints to the prompt via an external classifier. TabLLM makes use of parameter-efficient fine-tuning, such as IA3~\cite{liu2022few}, to adapt the LLM on limited training samples.

\vspace*{-0.12in}
\paragraph{Implementation details.}
\textsf{FeatLLM} adopts GPT-3.5 as its foundational LLM, but its architecture accommodates alternative LLM backbones (see Appendix~\ref{sec:backbones} for experiments using different models). LLM inference operates with a temperature of 0.5 and the top-p parameter is set to its standard value of 1 as specified by the API. The ensemble size and the number of extraction rules are set to 20 and 10, respectively. The influence of hyper-parameter choices is discussed in Figure~\ref{fig:hyperparameter2} and further elaborated in Appendix~\ref{sec:hyper-parameter}. For training the linear model, the Adam optimizer is employed with a learning rate of 0.01 across 200 epochs. Epoch selection utilizes $k$-fold cross-validation for best performance. When re-implementing baseline approaches, in-context learning is executed using GPT-3.5, and TabLLM leverages T0~\cite{sanh2022multitask}, following its original configuration. It is important to highlight that TabLLM only permits LLMs with publicly accessible checkpoints, which excludes GPT-3.5 from this setup. For conventional machine learning algorithms like LogReg, XGBoost, and RandomForest, parameter tuning is conducted through Grid-search combined with $k$-fold cross-validation. The value of $k$ is either 2 or 4, ensuring that every class has representation in the training data. Additional implementation details can be found in Appendix~\ref{sec:baselines}.

\vspace*{-0.12in}
\paragraph{Results.}
Tables~\ref{tab:main1} and \ref{tab:main2} provide a comparative analysis of model performances across multiple datasets. Each experimental setting was executed three times, each with unique random splits for the training and test sets; we report the mean and standard deviation of AUC across these runs. In nearly all circumstances, \textsf{FeatLLM} either outperforms other baselines or secures the second-best results. Notably, our proposed framework attains performance on par with traditional tabular learning methods that utilize 32 shots, even when restricted to just 4 shots (refer to Fig.~\ref{fig:compares}). While STUNT and TabPFN offered strong performance on some individual datasets, their cumulative improvements relative to standard machine learning models were limited. Further, LLM-based methods such as in-context learning, TABLET, and TabLLM proved ineffective in handling tabular datasets with high-dimensional feature spaces. In contrast, our framework, which leverages bagging and ensemble methodologies, remains robust and effective even when confronted with a large number of features. Additionally, while our approach to bagging and ensembling could be applied to other LLM-based models, it would result in a twofold increase in inference steps per sample, considerably escalating computational requirements and reducing practicality.

\subsection{Analysis \& Discussion}
\label{sec:analysis}
\paragraph{Ablation study.} To understand the influence of each principal component, we performed a series of ablation experiments, systematically removing: (1) \textsf{-Tuning}, which skips the weight optimization in the linear model, aggregating binary features directly for logit calculation; (2) \textsf{-Ensemble}, by removing the ensembling mechanism; (3) \textsf{-Description}, omitting feature descriptions; and (4) \textsf{-Reasoning}, bypassing the Step-1 reasoning instruction when constructing rules.

The outcomes, as detailed in Table~\ref{tab:ablation}, reflect average AUC changes across all datasets due to these ablations. Across the board, eliminating any of the core elements leads to diminished model performance. Notably, the positive impact of weight tuning becomes increasingly apparent with more available shots, suggesting that greater data enables more precise rule importance estimation, and thus makes tuning more advantageous. Conversely, the removal of feature descriptions and reasoning instructions leads to marked performance drops when shot counts are small, implying that leveraging the LLM’s inherent prior knowledge is most beneficial in such data-scarce regimes. Overall, our ensemble strategy yields consistent performance gains in every examined configuration.

\vspace*{-0.12in}
\paragraph{Training \& inference time comparison.} We evaluate the efficiency of training and inference times across several methods. The tests were conducted using the Adult dataset, with a training set size of 16 shots. A single A100 GPU served as the standard environment, except for TabLLM, which employed four A100 GPUs for model parallelism. Table~\ref{tab:cost} displays the comparative runtimes. \textsf{FeatLLM} demonstrates notably low inference times, on par with traditional baselines. Conversely, LLM-driven alternatives such as In-context learning, TABLET, and TabLLM are characterized by substantially longer inference times. For training time, \textsf{FeatLLM} maintains parity with other LLM-based strategies—it necessitates only 30 API calls, imposing minimal resource demands and allowing parallel execution.

\vspace*{-0.12in}
\paragraph{Quality of features generated by \textsf{FeatLLM}.}
To assess the utility of the rule-based features produced by \textsf{FeatLLM} in real-world tasks, we designed a straightforward analysis. For each dataset, we computed the mean AUROC between the generated features and the target labels, then determined the proportion of features achieving an AUROC greater than 0.5. Table~\ref{tab:quality} presents the findings. The results illustrate that most rules generated are pertinent to the target variable and furnish meaningful contributions to task performance (see ``Before tuning" column). Furthermore, when a linear model is tuned with the dataset, rules that lack significance (defined as having weights below a specified threshold) are systematically eliminated (see ``After tuning" column). The threshold, chosen as 0.1, was based on the distribution of learned weights.

\vspace*{-0.12in}
\paragraph{Impact of incorporating spurious correlations.} To evaluate how well different approaches handle extraneous spurious correlations, particularly under few-shot conditions, we performed an empirical study. The Heart dataset served as our base, into which we randomly integrated columns extracted from the Adult dataset—these added columns bear no relevance to the Heart task. We examined how the models' predictive performance responds as the quantity of unrelated columns increases. To ensure comparability across methods, we standardized the experiment to 8 labeled samples per class and did not include an unlabeled set, which consequently omits techniques like SCARF and STUNT from our assessment. 

Figure~\ref{fig:spurious} illustrates how performance shifts in response to spurious correlations. Among all evaluated models, \textsf{FeatLLM} exhibited the smallest performance drop when exposed to spurious features. This outcome can be attributed to our approach: when forming rules, it explicitly prioritizes the association between task objectives and feature sets according to prior domain knowledge, which discourages the inclusion of irrelevant features and leads to a more resilient model. In practice, we found that the fraction of rules derived from the noisy columns remained below 1-2\% of the total. On the other hand, it is also evident that \textsf{FeatLLM} is susceptible to erroneously capturing spurious associations or confounding causal links when additional columns from conceptually similar datasets are attached at random (refer to Appendix~\ref{sec:spurious-appendix} for further analysis).

\vspace*{-0.12in}
\paragraph{Hyper-parameter analysis}
\textsf{FeatLLM} incorporates two principal hyper-parameters: the ensemble count and the number of rules per class generated in each LLM query. To assess their effects, we performed empirical studies examining how varying these parameters influences system performance. Our findings indicated that increasing ensemble numbers leads to performance improvements; however, after about 20 ensembles, additional gains plateau (refer to Fig.~\ref{fig:appendix-hyperparameter1} in the Appendix). Regarding rule quantity, although variations did not yield statistically significant changes in outcomes, selecting 10 rules was deemed optimal as it strikes a balance between the length of prompts and performance efficacy (see Fig.~\ref{fig:hyperparameter2}). We hypothesize that this result stems from the inclusion of more rules expanding the input space and thereby enhancing informational content, while also heightening the risk of overfitting, particularly under low-shot conditions.

\section{Conclusion}
We introduce \textsf{FeatLLM}, advancing state-of-the-art few-shot tabular learning with LLMs. Rather than applying LLMs solely for direct sample-level prediction, \textsf{FeatLLM} leverages the models to generate predictive criteria analogous to feature engineering. Via rule generation paired with an ensemble strategy using bagging, \textsf{FeatLLM} demands minimal inference effort, operates fully through API calls without necessitating any model training, and is not limited by the dimensionality of the tabular data. The method delivers high accuracy, offering a robust, efficient solution for deploying LLMs on practical tabular datasets.

\textsf{FeatLLM} is optimized specifically for the low-shot learning scenario. Going forward, we aim to broaden its scope to facilitate new feature discovery in settings with larger sample sizes, experiment with feature types beyond rule-based ones, and increase interpretability, thereby unlocking its application across diverse tasks.

\section{Broader Impact}
The \textsf{FeatLLM} framework is designed to harness the prior knowledge encoded in LLMs for tabular learning, surpassing what traditional supervised approaches typically achieve. By prioritizing data-efficient learning, our research addresses the challenge of overfitting in contexts with sparse training samples, supplementing these with rich prior information. This approach proves particularly valuable for domains such as finance or healthcare, where label acquisition is expensive or logistically difficult, and domain-relevant prior knowledge within LLMs can be leveraged effectively due to pretraining on pertinent textual data. A critical direction for future research will be thorough investigations into the effects of societal biases and inaccuracies embedded in LLM priors, as these may exert substantial influence—sometimes outweighing empirical data—on model predictions, which is key for safe and equitable widespread deployment.

\end{document}